\chapter{Scope and Method}
\label{ch:scope}

\section{From a demonstration to an engineering objective}

The first cylinder answered the original feasibility question: reinforcement learning could
produce a basic landing behaviour. It could not answer whether that result would survive a
change of vehicle, reward, actuator or simulator version. The engineering objective therefore
shifted from producing one successful policy to producing a desktop simulator in which an
experiment could be configured, repeated and inspected.

In the resulting workflow, a user defines the vehicle, environment, task, reward, learning run
and telemetry before launching one or more independent simulation areas. The same application
supports PPO training, ONNX inference, deterministic evaluation and component-level tests. A
completed run should retain the configuration that produced it, and a policy whose observation
or action interface no longer matches the vehicle should be rejected before inference begins.

The simulator eventually grew beyond what could be evaluated within one bachelor thesis. The
experimental scope was therefore reduced to a sequence of three increasingly demanding tasks:
holding a fixed hover, tracking relocated hover targets, and landing on four physical legs.
Alternative vehicle configurations, wind and injected actuator faults remain implemented
capabilities, but they are not presented as validated results. External-arm capture is treated
as a separate future task rather than being added to an already broad landing study.

\section{Development method}

Development followed the same loop throughout the project:

\begin{enumerate}
  \item add or modify one physical, actuator or task component;
  \item test its direction, limits and feasibility independently of learning;
  \item train a policy and inspect trajectories, episode outcomes and telemetry;
  \item identify the simulator defect, reward loophole or control limitation revealed by the
  new behaviour;
  \item move configuration and responsibility into a clearer module when the change exposed
  unwanted coupling; and
  \item preserve the resulting configuration as a new run or immutable revision before
  continuing.
\end{enumerate}

This method emerged because training alone was a poor debugging instrument. A bad learning
curve could mean that the task was too difficult, but it could also mean that drag pointed in
the wrong direction or that the requested spawn velocity was never applied. I therefore added
what the interface calls \emph{hardware tests}: component-level experiments inside Unity that
isolate the simulated hardware from the policy. No physical rocket hardware was built or
tested. These checks answer questions such as whether drag opposes air-relative motion, whether
an engine respects its spool rate, and whether the available thrust can arrest the selected
descent.

The important part of the loop is that an unexpected policy is treated as evidence. It is first
reproduced and measured, then used to decide what the next run should change. Chapter~\ref{ch:development}
applies this method to the L1--L11 landing campaign.

\section{Separating capability from evidence}

As the simulator gained features, it became easy to confuse an available checkbox with an
evaluated contribution. To prevent that, this thesis uses three evidence classes. A
\emph{completed evaluation} executes its predefined episode count and fixed seed without manual
interruption. A \emph{preliminary diagnostic} can reveal useful behaviour but is too small or
incomplete for a reliability claim. A \emph{capability} exists in the code or UI but has not yet
been evaluated under the thesis protocol.

The \texttt{HoverFixed}, moving-target hover and L11 landing policies have completed their
planned evaluation episode counts. L10 is retained as a completed landing baseline because
comparing it with L11 shows what changed after the actuator and reward redesign. Fixed hover is
reported through continuous state and control metrics because its evaluator still has no binary
success rule; completing 250 horizons is not silently relabelled as 100\% task success. Wind,
faults and vehicle-configuration comparisons remain capabilities. These labels follow each
result into Chapters~\ref{ch:evaluation}--\ref{ch:limitations}.

\section{Reproducibility boundary}

Each reported run is identified by its run ID, session revision, environment and trainer seed,
training configuration, exported model and evaluation summary. This record is the final step of
the development loop: it connects a plotted result back to the exact experiment that produced
it. The public repository URL and commit cannot be filled in truthfully until the working tree
is frozen, because the current Git HEAD does not include uncommitted simulator changes.

\begin{thesistodo}
  Freeze the submission revision, publish the repository, and replace the repository and commit
  placeholders in Appendix~\ref{app:reproducibility}. Do this only after the working tree is
  stable, so that the published identity matches the submitted implementation.
\end{thesistodo}
